\section{Introduction}

    \subsection{Background and Motivation}
        Robotic arm control in the presence of unknown disturbances remains a fundamental challenge in modern control theory and robotics. 
        Disturbances such as friction, payload variations, unmodeled dynamics, and external forces can severely degrade tracking performance 
        and reduce system reliability. Traditional PD controllers, while simple and effective in ideal conditions, provide limited robustness 
        against model uncertainties and disturbances.

        Disturbance Observer-Based (DOB) control has emerged as a promising approach to address these challenges. By estimating unknown 
        disturbances in real-time, observers enable feedback compensation without requiring explicit system models. This work implements 
        three observer variants---Q-filter Disturbance Observer (DOB), Momentum Observer, and Force Observer---to evaluate their effectiveness 
        in trajectory tracking and disturbance rejection tasks.

    \subsection{Simulation Platform}
        The simulation is built using:\vspace{-4pt}
        \begin{itemize}
            \item \textbf{MuJoCo 3.7.0}: Physics engine providing accurate dynamics modeling
            \item \textbf{AgileX Piper}: 6-DOF collaborative manipulator with joint torque actuation
            \item \textbf{C++17}: Core simulation loop with modular control architecture
            \item \textbf{GLFW + OpenGL}: Real-time 3D visualization with trajectory monitoring
        \end{itemize}

    \subsection{Research Objectives}
        \begin{enumerate}
            \item Validate three observer implementations (DOB, Momentum, Force) in a high-fidelity simulator
            \item Quantify tracking error under matched and mismatched disturbances
            \item Compare computational cost vs. tracking performance across observer types
            \item Demonstrate gravity compensation and friction estimation capabilities
            \item Provide reproducible experimental framework for future research
        \end{enumerate}

    \subsection{Report Organization}
        \begin{description}
            \item[Chapter 2] System description, control architecture, and observer theory
            \item[Chapter 3] Simulation setup, experimental parameters, and disturbance models
            \item[Chapter 4] Joint-space tracking results and per-joint analysis
            \item[Chapter 5] End-effector 3D trajectory and workspace analysis
            \item[Chapter 6] Conclusions, performance comparisons, and future work
        \end{description}